Establecido el lema, en el caso al que nos hemos reducido al final
de II), $B'$ es \emph{radicial} sobre $A'$ (\textbf{I}, 3.5.7) y, por tanto, es también
un anillo \emph{local}; además, $B'$ es entero sobre $A'$; así, si
$\mathfrak q'$ es el único ideal primo de $B'$ situado sobre
$\mathfrak p'$, $\mathfrak q'B'_{\mathfrak p'}$ es el único ideal
maximal de $B'_{\mathfrak p'}$ (Bourbaki, \emph{Alg. comm.}, cap. V, § 2,
n\textsuperscript{o} 1, prop. 1), y en consecuencia
$B'_{\mathfrak q'}=B'_{\mathfrak p'}=B\otimes_A
A'_{\mathfrak p'}$. Demostraremos primero que
$B'_{\mathfrak p'}$ es un anillo \emph{noetheriano} y \emph{normal}. Puesto que $B$
contiene a $A$ y está contenido en $K$, y que $A'_{\mathfrak p'}$ es
un $A$-módulo plano, $B'_{\mathfrak p'}$ contiene a
$A'_{\mathfrak p'}$ y está contenido en $K'_{\mathfrak p'}$ y, por
tanto, en el producto $L_1$ de los $L'_j$ para los que
$\mathfrak q'_j\subset\mathfrak p'$.
